\section{Abstract}
This document presents the development of the project "Web System for Document Management and QR Code Generation for PJETAM", carried out at the Poder Judicial del Estado de Tamaulipas, located at Boulevard Praxedis Balboa 2207, between López Velarde and Díaz Mirón streets, Colonia Miguel Hidalgo, ZIP code 87090, in Ciudad Victoria, Tamaulipas.
The main objective of the project was to develop a web application that would allow staff to convert, store and manage documents and images, generating QR codes linked to each file to facilitate their referencing. The application was integrated with existing desktop management systems through authentication via an encrypted parameter in the URL.

The system was built using the MVC architecture on ASP.NET Core 10 with the C\# programming language, SQL Server for the database managed through Entity Framework Core, AWS S3 for cloud file storage, LibreOffice in headless mode for the PDF conversion service and the QRCoder library for QR code generation. For the development of user interfaces, Razor views (.cshtml) combined with Bootstrap 5 were used. Knowledge acquired during higher education in relational databases, cloud storage, object-oriented programming and web development made it possible to integrate different technologies to build a complete and functional system.

Upon completion of the project, the File Conversion and Upload module and the Document Management module were implemented, centralizing file searching, visualization, download, deletion and QR code generation. Thus, the system met the established requirements, providing a single application with all document management processes, eliminating the need for external tools and inefficient workflows.